\chapter*{图片版权说明}

本书中的图表和示意图皆为原创；依据已发表数据绘制的地图与图表，在图注或本列表中注明来源。照片依各自的自由许可协议使用，谨向作者致谢；每张照片的完整来源链接和许可协议网址，列于本书代码库中的 \texttt{images/book1/CREDITS.md}。

\bigskip
与照片和原创图表并列，部分图为 AI 生成的示意图，由 OpenAI 图像生成功能依本书编辑撰写的提示词生成，并在收录前经过科学准确性审核。它们只呈现通用场景，从不呈现真实地点、真实人物或真实标本。每幅生成图像的完整提示词记录于代码库中的 \texttt{images/book1/ai/PROMPTS.md}。

\bigskip
\noindent 照片与数据，按章列出：
\begin{itemize}\small
  \item 第~\ref{ch:g2:water-on-earth}~章（地球上的水）：阿波罗 17 号于 1972 年 12 月 7 日所拍的地球照片（“蓝色弹珠”）；NASA，公有领域（经维基共享资源）。100 个方块依 NOAA ETOPO1 地形数据计算；100 个水桶取自美国地质调查局水科学学堂的表格（引自 Shiklomanov，1993）。
  \item 第~\ref{ch:g3:stones-pebbles-sand}~章（石块、卵石与沙）：一堵干砌花岗岩石墙；照片由 Tony Atkin 拍摄，Geograph，CC BY-SA 2.0（经维基共享资源）。瑞士的石板屋顶；照片由 Michael C. Rygel 拍摄，CC BY-SA 3.0（经维基共享资源）。
  \item 第~\ref{ch:g3:continents-and-oceans}~章（地球仪）：两幅地球仪与世界地图依 Natural Earth 海岸线绘制（公有领域）。
  \item 第~\ref{ch:g4:clouds-rain-snow}~章（云、雨与雪）：堪萨斯上空的积云（S.~Corfidi）、卷云片（R.~F.~Kresge）、山峰周围的层云、亚利桑那沙漠上空的雷暴（B.~Hooke 藏品），以及一颗巨大的冰雹；均来自 NOAA 图片库，公有领域（经维基共享资源）。
  \item 第~\ref{ch:g4:volcanoes-earthquakes-first-look}~章（火山与地震初识）：基拉韦厄火山熔岩入海，美国地质调查局夏威夷火山观测站；1980 年 5 月 18 日的圣海伦斯火山，照片由 Austin Post 拍摄，美国地质调查局；1989 年洛马普列塔地震将一座房屋推离地基，照片由 J.~K.~Nakata 拍摄，美国地质调查局；均为公有领域（经维基共享资源）。地图采用史密森尼学会全球火山活动计划的全新世火山名录，以及美国地质调查局的地震目录。
  \item 第~\ref{ch:g4:world-physical-map}~章（世界自然地图）：国际空间站拍摄的喜马拉雅山与亚马孙河，NASA，公有领域（经维基共享资源）。自然地图依 NOAA ETOPO1 地形数据绘制，河流、湖泊与海岸线取自 Natural Earth（公有领域）。
  \item 第~\ref{ch:g5:water-shapes-the-land}~章（水塑造大地）：国际空间站拍摄的马莫雷河（ISS007-E-10797），以及 Terra 卫星上 MODIS 仪器所见的尼罗河三角洲；NASA，公有领域（经维基共享资源）。
  \item 第~\ref{ch:g5:weather-and-climate}~章（天气与气候）：索诺兰沙漠，美国国家公园管理局萨瓜罗国家公园，公有领域（经维基共享资源）。气候图依 NOAA NCEI 美国气候常态值 1991--2020 绘制（公有领域）。
  \item 第~\ref{ch:g5:natural-hazards}~章（自然灾害）：国际空间站上 Ed Lu 拍摄的伊莎贝尔飓风之眼，NASA，公有领域；2014 年华盛顿州的滑坡，Mark Reid，美国地质调查局，公有领域（均经维基共享资源）。
  \item 第~\ref{ch:g6:earth-a-planet}~章（地球，一颗行星）：水星 10 号拍摄的金星（NASA/JPL-Caltech）；阿波罗 17 号拍摄的地球（NASA）；阿波罗 8 号拍摄的《地出》（NASA，Bill Anders）；哈勃太空望远镜拍摄的火星与木星（NASA、ESA、哈勃遗产团队及 A.~Simon）；均为公有领域（经维基共享资源）。1990 年旅行者 1 号所摄的《淡蓝色圆点》，2020 年版（NASA/JPL-Caltech，PIA23645，公有领域）。行星数据取自 NASA 行星概况表。
  \item 第~\ref{ch:g6:latitude-longitude}~章（纬度与经度）：格林尼治国家海事博物馆所藏哈里森 H4 航海天文钟，照片由 Colonel Warden 拍摄，CC BY-SA 3.0（经维基共享资源）。海岸线与城市取自 Natural Earth（公有领域）。
  \item 第~\ref{ch:g6:rocks-and-minerals}~章（岩石与矿物）：花岗岩、玄武岩、含化石石灰岩、砂岩、大理石、页岩、一块石英晶体与一片白云母薄片，照片由 James St.~John 拍摄，CC BY 2.0；一块板岩屋顶，照片由 Øyvind Holmstad 拍摄，CC0（均经维基共享资源）。
  \item 第~\ref{ch:g6:inside-the-earth}~章（地球内部）：1932 年的英格·莱曼，照片由 Even Neuhaus 拍摄，丹麦皇家图书馆，公有领域（经维基共享资源）。各层的密度与波速依初步参考地球模型绘制（A.~M. Dziewonski 与 D.~L. Anderson，1981）。
  \item 第~\ref{ch:g6:water-underground}~章（地下水）：密苏里州的大泉，照片由 Johnston9494 拍摄，CC BY-SA 4.0（经维基共享资源）。孔隙度引自 R.~C. Heath，《地下水水文学基础》，美国地质调查局（1983）；淡水占比取自美国地质调查局水科学学堂。
  \item 第~\ref{ch:g7:plate-tectonics}~章（板块构造）：重新拼合的南方大陆，并标示出犬齿兽、Lystrosaurus、中龙与舌羊齿的化石带；美国地质调查局，《动态地球》，公有领域。
  \item 第~\ref{ch:g7:plate-tectonics}~章：横跨卡里佐平原的圣安德烈亚斯断层，照片由 Robert E. Wallace 拍摄，美国地质调查局，公有领域。
  \item 第~\ref{ch:g7:plate-tectonics}~章：板块边界依 P.~Bird（2003），GeoJSON 由 H.~Ahlenius 制作（ODC-By 1.0）；海岸线取自 Natural Earth（公有领域）；地震取自美国地质调查局目录；火山取自史密森尼学会全球火山活动计划。
  \item 第~\ref{ch:g7:volcanoes}~章（火山）：1980 年 5 月 17 日与 9 月 10 日从约翰斯顿岭拍摄的圣海伦斯火山，照片由 Harry Glicken 拍摄，美国地质调查局；1983 年基拉韦厄火山的熔岩喷泉，以及 1985 年从空中拍摄的莫纳罗亚火山，照片由 J.~D. Griggs 拍摄，美国地质调查局；1991 年 6 月 12 日皮纳图博火山的喷发柱，照片由 Dave Harlow 拍摄，美国地质调查局；均为公有领域（经维基共享资源）。地图与计数采用史密森尼学会全球火山活动计划的全新世火山名录；板块边界依 P.~Bird（2003）。
  \item 第~\ref{ch:g7:earthquakes}~章（地震）：2004 年 12 月 26 日海啸后苏门答腊海岸附近的一座村庄，美国海军照片，Philip A. McDaniel 拍摄；1983 年博拉峰地震的断层陡崖，照片由 R.~E. Wallace 拍摄，美国地质调查局；1959 年赫布根湖边塌入湖中的公路，照片由 I.~J. Witkind 拍摄，美国地质调查局；均为公有领域（经维基共享资源）。计数与深度取自美国地质调查局地震目录；烈度表节选自美国地质调查局的修正麦卡利烈度表。
  \item 第~\ref{ch:g7:sediments-and-fossils}~章（沉积物与化石）：隐士小憩处俯瞰的大峡谷，国家公园管理局照片，Kristen M. Caldon 拍摄，公有领域；Capitol Reef 国家公园中的一块菊石化石、一块三叶虫化石与波痕，照片由 James St.~John 拍摄，CC BY 2.0；恐龙岭的恐龙足迹，照片由 Dexter Perkins 拍摄，CC0（均经维基共享资源）。
  \item 第~\ref{ch:g7:geological-time}~章（地球的年龄）：詹姆斯·赫顿的肖像，Henry Raeburn 绘，公有领域（经维基共享资源）。地质年代表的年龄与颜色取自国际年代地层表（国际地层委员会）。
  \item 第~\ref{ch:g8:the-atmosphere}~章（大气层）：在堪萨斯州托皮卡释放探空气球，照片由 Stephen Corfidi 拍摄，NOAA/国家气象局；国际空间站拍摄的日落时分的地球边缘，NASA，国际空间站第 23 期考察；均为公有领域（经维基共享资源）。廓线依《美国标准大气（1976）》计算（NOAA、NASA 与美国空军）。
  \item 第~\ref{ch:g8:winds-global-circulation}~章（风与全球环流）：2022 年 7 月 26 日 GOES-16 卫星所见的美洲，NOAA，公有领域（经维基共享资源）。信风图的海岸线取自 Natural Earth（公有领域）。
  \item 第~\ref{ch:g8:fronts-and-weather-maps}~章（锋面与天气图）：Terra 卫星所见冰岛附近的一个低压区，NASA/GSFC，MODIS 快速响应小组，Jacques Descloitres 拍摄，公有领域（经维基共享资源）。气团图的海岸线取自 Natural Earth（公有领域）。
  \item 第~\ref{ch:g8:the-oceans}~章（海洋）：1872 年 HMS \emph{Challenger} 号准备测深，取自挑战者号考察报告的雕版画，转载于 W.~A. Herdman，《海洋学的奠基人》（1923）；Aqua 卫星上 MODIS 仪器所见的墨西哥湾流，NASA 地球观测站；均为公有领域（经维基共享资源）。洋底剖面依 NOAA ETOPO1 地形数据绘制；环流图与浮标图的海岸线取自 Natural Earth（公有领域）。
  \item 第~\ref{ch:g8:ice-on-earth}~章（地球上的冰）：Pedersen 冰川在 1990 年与 2016 年的景象，照片由 M.~Tetreau 与 D.~Kurtz 拍摄，美国国家公园管理局；南极半岛外某冰架的边缘，NASA，IceBridge 行动（M.-J. Viñas）；均为公有领域（经维基共享资源）。海冰曲线采用 NSIDC 海冰指数第 4 版（F.~Fetterer 等，2025，doi:10.7265/a98x-0f50）。
  \item 第~\ref{ch:g9:climate-past-and-present}~章（地球气候的过去与现在）：格陵兰 GISP2 冰芯的一段，NOAA 古气候学，公有领域（经维基共享资源）。曲线依以下数据绘制：NASA GISTEMP 第 4 版、莫纳罗亚二氧化碳记录（NOAA 全球监测实验室与斯克里普斯海洋研究所）、EPICA 圆顶 C 站温度（J.~Jouzel 等，2007）、南极冰芯二氧化碳综合记录（B.~Bereiter 等，2015）以及 NOAA 卫星测高实验室的海平面数据；评估表述引自 IPCC 第六次评估报告第一工作组《决策者摘要》（2021）。
  \item 第~\ref{ch:g9:natural-risks}~章（与自然风险共处）：内瓦多德尔鲁伊斯火山的危险区图，取自 T.~L. Wright 与 T.~C. Pierson，《与火山共存》，美国地质调查局通报 1073（1992）；美国海岸警卫队 \emph{William Hart} 号船员检查的海啸预警浮标，美国海岸警卫队，2021；均为公有领域（经维基共享资源）。
  \item 第~\ref{ch:g9:earth-resources}~章（地球的资源）：ASTER 所见的丘基卡马塔矿区，NASA/METI/AIST/日本空间系统及美日 ASTER 科学团队，2007，公有领域；俄亥俄州贝尔蒙县公路切坡中的一层煤层，照片由 James St.~John 拍摄，CC BY 2.0（均经维基共享资源）。富集系数取自 D.~Perkins 等，《矿物学》，以及 S.~Earle，《环境地质学》。
  \item 第~\ref{ch:g9:landforms}~章（地貌）：七姐妹峰，照片由 Arild Vågen 拍摄，2017，CC BY-SA 4.0；巴西南部海岸的新月形沙丘，2017 年国际空间站宇航员拍摄，公有领域；卡尔斯巴德洞穴中的钟乳石，国家公园管理局照片，Daniel Leifheit 拍摄，2018，公有领域（均经维基共享资源）。
  \item 第~\ref{ch:g9:maps-of-the-earth}~章（整个地球的地图）：2000 年陆地卫星影像所示的勒拿河三角洲，USGS EROS 数据中心，经 NASA Visible Earth，公有领域；一颗 GPS 卫星，NASA 艺术家构想图，公有领域（均经维基共享资源）。世界地图上的海岸线取自 Natural Earth（公有领域）。
\end{itemize}
